\documentclass[11pt,a4paper]{article}
\usepackage[utf8]{inputenc}
\usepackage[T1]{fontenc}
\usepackage[french]{babel}
\usepackage{amsmath,amssymb}
\usepackage{geometry}
\usepackage{booktabs}
\usepackage{hyperref}
\usepackage[table]{xcolor}
\usepackage{colortbl}
\usepackage{listings}
\usepackage{fancyhdr}

\definecolor{sectioncolor}{RGB}{0,51,102}
\definecolor{subsectioncolor}{RGB}{0,102,102}
\definecolor{codebg}{RGB}{245,245,250}
\definecolor{tableheader}{RGB}{0,51,102}
\definecolor{newmodule}{RGB}{0,128,0}
\definecolor{revised}{RGB}{153,51,0}

\geometry{margin=2.5cm}
\setlength{\parindent}{0pt}
\setlength{\parskip}{0.6em}

\lstset{
  basicstyle=\ttfamily\small,
  breaklines=true,
  backgroundcolor=\color{codebg},
  frame=single,
  framerule=0.5pt,
}

\newcommand{\newmod}[1]{\textcolor{newmodule}{\textbf{#1}}}
\newcommand{\revised}[1]{\textcolor{revised}{\textbf{#1}}}

\hypersetup{colorlinks=true, linkcolor=sectioncolor, urlcolor=subsectioncolor}

\begin{document}

\begin{titlepage}
\centering
\vspace*{2cm}
{\Huge\bfseries\color{sectioncolor} Rapport Technique}\\[0.5em]
{\LARGE Chatbot Juridique Multilingue}\\[0.3em]
{\large Architecture Révisée et Améliorations Proposées}\\[2cm]
\begin{tabular}{ll}
\textbf{Projet :} & Chatbot Juridique Multilingue avec RAG \\
\textbf{Équipe :} & Bettayeb Mohamed Aimen, Amina, Yanis, Loubna \\
\textbf{Encadrante :} & Mme Berkani \\
\textbf{Date :} & Mars 2026 \\
\textbf{Version :} & 2.0 --- Architecture Révisée \\
\end{tabular}
\vfill
\end{titlepage}

\tableofcontents
\newpage

\section{Résumé Exécutif}

Ce rapport présente l'architecture complète et révisée du chatbot juridique multilingue développé dans le cadre du projet de fin d'études. Suite aux remarques formulées lors de la dernière séance concernant la nature arbitraire de certaines formules mathématiques utilisées dans le module de classification d'intention, nous avons procédé à une analyse critique de l'ensemble de l'architecture existante. Ce document décrit les modifications apportées, les justifications théoriques de chaque décision de conception, et les améliorations attendues.

\section{Contexte et Problématique}

\subsection{Motivation du Projet}

L'accès aux régulations légales et administratives au sein des institutions académiques représente un défi majeur pour les étudiants et les chercheurs. Les textes juridiques sont rédigés dans un langage complexe, distribués sur de nombreux documents hétérogènes, et disponibles en plusieurs langues (arabe, français, anglais).

\textbf{Exemple concret :} Un doctorant dont le directeur de thèse décède en cours de parcours doit suivre des procédures administratives précises définies par les textes réglementaires universitaires. Sans outil adapté, il doit manuellement parcourir de nombreux documents pour identifier les démarches requises.

Notre système propose une solution conversationnelle permettant aux utilisateurs de poser des questions en langage naturel dans leur langue maternelle et d'obtenir des réponses fondées et vérifiées.

\subsection{Problème Central}

Les modèles de langage génératifs peuvent produire des informations juridiques incorrectes avec une haute confiance apparente. Dans un contexte juridique, une réponse erronée peut avoir des conséquences réelles. L'architecture Retrieval-Augmented Generation (RAG) est adoptée : les réponses sont ancrées dans des documents externes vérifiés.

\subsection{Critique de l'Architecture Précédente}

Suite à la remarque de Mme Berkani lors de la dernière séance, nous avons identifié que la principale faiblesse résidait dans le module de classification d'intention. Ce module utilisait une formule heuristique :
\begin{equation}
\text{score}_i = \text{base}_i + 0.03 \times \min(m_i - 1, 3)
\end{equation}
Les constantes choisies (0.03, 0.85, 0.15) n'avaient pas été validées empiriquement. De plus, les patterns regex étaient principalement rédigés en anglais et français, créant une lacune pour le traitement des requêtes en arabe pur.

\section{Architecture Globale du Pipeline}

Le pipeline révisé comprend : Détection de langue (Unicode + langdetect) $\to$ Résolution de contexte conversationnel $\to$ Classification d'intention (LLM Zero-Shot) $\to$ Routage $\to$ Représentation vectorielle (BGE-M3) $\to$ Retrieval hybride (Dense + BM25 + RRF) $\to$ Déduplication (Jaccard) $\to$ Reranking $\to$ Construction du contexte $\to$ Génération LLM $\to$ Vérification de fidélité $\to$ Réponse finale.

\section{Module 1 --- Détection de Langue}

\subsection{Technique 1 --- Analyse de Script Unicode}

Les caractères arabes sont détectés par une expression régulière couvrant les plages Unicode arabes. Le ratio arabe :
\begin{equation}
R_{\text{arabic}} = \frac{N_{\text{ar}}}{N_{\text{alpha}}} \geq 0.30 \Longrightarrow \text{lang} = \text{ar}
\end{equation}
Le seuil 0.30 offre le compromis optimal entre robustesse au code-mixing et précision, conformément aux pratiques documentées (langdetect, pycld2).

\subsection{Technique 2 --- Détection Statistique}

Si le texte n'est pas arabe, la bibliothèque \texttt{langdetect} est utilisée. Formule (Maximum Likelihood) :
\begin{equation}
L^* = \arg\max_{L} P(L \mid T) = \arg\max_{L} P(T \mid L) \cdot P(L)
\end{equation}

\section{Module 2 --- Classification d'Intention (Révisé)}

\subsection{Problèmes avec l'Ancienne Approche}

\textbf{Problème 1 --- Lacune linguistique :} Les patterns regex étaient en anglais et français. Une requête en arabe ne correspondait à aucun pattern et était classifiée par défaut comme \texttt{general\_knowledge}.

\textbf{Problème 2 --- Constantes sans validation :} Le coefficient 0.03, les scores de base, le seuil 0.15 ont été choisis par intuition sans validation sur données réelles.

\textbf{Problème 3 --- Maintenance :} Chaque nouvelle formulation nécessite l'ajout manuel d'un pattern regex.

\subsection{Nouvelle Approche --- LLM Zero-Shot}

La requête est envoyée au LLM avec un prompt structuré décrivant les intents : \texttt{legal\_query}, \texttt{document\_query}, \texttt{general\_knowledge}. Le LLM retourne l'intent approprié en s'appuyant sur sa compréhension sémantique, indépendamment de la langue.

\subsection{Résolution de Contexte Conversationnel}

Avant la classification, une étape de réécriture de requête pour les conversations multi-tours : le système réécrit la question de suivi sous forme d'une question complète et autonome incluant le contexte de l'historique.

\section{Module 3 --- Routage des Requêtes}

\begin{center}
\begin{tabular}{ll}
\toprule
\rowcolor{tableheader}\textcolor{white}{\textbf{Intent}} & \textcolor{white}{\textbf{Source}} \\
\midrule
\texttt{legal\_query} & Qdrant --- \texttt{legal\_documents} \\
\texttt{document\_query} & Qdrant --- \texttt{document\_chunks} \\
\texttt{general\_knowledge} & LLM directement \\
\bottomrule
\end{tabular}
\end{center}

\section{Module 4 --- Représentation Vectorielle (Révisé)}

\subsection{Changement de Modèle}

\textbf{Ancien :} \texttt{paraphrase-multilingual-mpnet-base-v2} (768 dimensions, support arabe limité).

\textbf{Nouveau recommandé :} \texttt{BAAI/bge-m3} ou \texttt{intfloat/multilingual-e5-large}. BGE-M3 (2024) : entraîné sur un corpus multilingue plus large, optimisé pour le retrieval, génère vecteurs denses et sparses, performances supérieures sur MIRACL incluant l'arabe.

\subsection{Similarité Cosinus}

\begin{equation}
\cos(q, d) = \frac{q \cdot d}{\|q\| \|d\|} = \frac{\sum q_i d_i}{\sqrt{\sum q_i^2} \cdot \sqrt{\sum d_i^2}}
\end{equation}
Invariance à la norme ; borné dans $[0,1]$ ; standard IR (Salton \& McGill, 1983).

\section{Module 5 --- Retrieval Hybride Réel (Révisé)}

\subsection{Problème avec l'Ancienne Approche}

L'architecture précédente était une recherche sémantique avec pondération manuelle : $s_{\text{final}}(q,d,s) = \cos(q,d) \times w_s$ avec $w_s = 1.05$. Aucune composante lexicale ; poids arbitraire.

\subsection{Nouvelle Approche}

\textbf{Composante 1 --- Dense :} Qdrant avec embeddings BGE-M3.

\textbf{Composante 2 --- Sparse :} BM25 sur les tokens.
\begin{equation}
\text{BM25}(q,d) = \sum_{t \in q} \text{IDF}(t) \cdot \frac{f(t,d) \cdot (k_1+1)}{f(t,d) + k_1 \cdot (1 - b + b \cdot |d|/\text{avgdl})}
\end{equation}
Avec $k_1 = 1.5$, $b = 0.75$.

\subsection{Fusion RRF}

\begin{equation}
\text{RRF\_score}(d) = \frac{1}{k + \text{rank}_{\text{dense}}(d)} + \frac{1}{k + \text{rank}_{\text{sparse}}(d)}
\end{equation}
Avec $k = 60$ (Cormack et al., 2009). Aucun paramètre arbitraire ; insensible aux échelles ; robustesse.

\section{Module 6 --- Déduplication}

\begin{equation}
J(A, B) = \frac{|A \cap B|}{|A \cup B|}, \qquad J(A,B) \geq 0.85 \Longrightarrow \text{duplicat}
\end{equation}
Seuil 0.85 : pratique établie (Manku et al., 2007).

\section{Module 7 --- Reranking}

\begin{equation}
\text{score}_{\text{rerank}}(q, d) = \frac{q \cdot d}{\|q\| \|d\| + \varepsilon}, \quad \varepsilon = 10^{-9}
\end{equation}
Recalcule la similarité cosinus sur les candidats restants. Top-5 chunks retenus.

\section{Module 8 --- Construction du Contexte (Révisé)}

\subsection{Balanced Selection}

$\forall f \in \mathcal{F}$, $|\mathcal{R}_f| \geq 2$ (MIN\_PER\_DOC = 2).

\subsection{Chunking Structure-Aware}

Priorité 1 : frontières d'articles (Article X, Art. X, Al-Mâda X). Priorité 2 : paragraphes. Priorité 3 : phrases. Priorité 4 : tokens (512, overlap 64) en fallback. CHUNK\_SIZE = 512, CHUNK\_OVERLAP = 64 (Lewis et al., 2020).

\section{Module 9 --- Génération LLM}

Le prompt intègre : contexte récupéré, requête (réécrite si multi-tours), historique, instructions de citation. Estimation des tokens : $\widehat{T}(x) = \max(1, \lfloor |x|/4 \rfloor)$.

\section{Module 10 --- Vérification de Fidélité (Nouveau)}

Après génération, le LLM vérifie si chaque affirmation est supportée par le contexte. Répond \texttt{faithful} ou \texttt{not\_faithful}. Si \texttt{not\_faithful} : réponse de fallback sécurisée recommandant la consultation des textes ou du service administratif.

\section{Tableau Récapitulatif des Changements}

\begin{center}
\small
\begin{tabular}{p{2.2cm}p{2.2cm}p{2.2cm}p{2.8cm}}
\toprule
\textbf{Composant} & \textbf{Précédent} & \textbf{Révisé} & \textbf{Raison} \\
\midrule
Classification & Regex + heuristique & LLM Zero-Shot & Agnostique langue \\
Retrieval & Sémantique pondérée & Dense + BM25 + RRF & Références exactes \\
Embedding & mpnet-base-v2 & BGE-M3 & Support arabe \\
Chunking & Fixe 512 & Hiérarchique & Frontières articles \\
Post-génération & Aucune & Vérification fidélité & Anti-hallucination \\
Conversations & Historique brut & Réécriture requête & Références pronominales \\
\midrule
Cosinus, Jaccard, Top-K=5, MIN\_PER\_DOC=2, BERTScore, P@k/R@k/MRR & Conservés \\
\bottomrule
\end{tabular}
\end{center}

\section{Méthodologie d'Évaluation}

\textbf{Retrieval :} $P@k$, $R@k$, $\text{MRR} = (1/Q) \sum_{i=1}^{Q} 1/\text{rank}_i$.

\textbf{Génération :} BERTScore (similarité sémantique vs BLEU/ROUGE lexical).

\textbf{Humaine :} Exactitude, complétude, clarté, citation des sources.

\textbf{Vérification fidélité :} Taux de détection hallucinations, faux positifs, fallback.

\section{Références Bibliographiques}

\begin{enumerate}
\item Salton, G. \& McGill, M.J. (1983). \textit{Introduction to Modern Information Retrieval}. McGraw-Hill.
\item Lewis, P. et al. (2020). ``Retrieval-Augmented Generation for Knowledge-Intensive NLP Tasks.'' NeurIPS.
\item Manku, G.S., Jain, A., \& Das Sarma, A. (2007). ``Detecting Near-Duplicates for Web Crawling.'' WWW.
\item Reimers, N. \& Gurevych, I. (2019). ``Sentence-BERT.'' EMNLP.
\item Zhang, T. et al. (2019). ``BERTScore.'' ICLR 2020.
\item Bruch, S., Gai, A., \& Ingber, A. (2022). ``Fusion Functions for Hybrid Retrieval.'' arXiv:2210.11934.
\item Nogueira, R. \& Cho, K. (2019). ``Passage Re-ranking with BERT.'' arXiv:1901.04085.
\item Cormack, G.V., Clarke, C.L.A., \& Buettcher, S. (2009). ``Reciprocal Rank Fusion.'' SIGIR.
\item Chen, J. et al. (2024). ``BGE M3-Embedding.'' arXiv:2402.03216.
\item Robertson, S. \& Zaragoza, H. (2009). ``BM25 and Beyond.'' Foundations and Trends in IR.
\end{enumerate}

\section*{Annexe --- Constantes}

\begin{center}
\begin{tabular}{lll}
\toprule
\textbf{Constante} & \textbf{Valeur} & \textbf{Source} \\
\midrule
Seuil arabe & 0.30 & Code-mixing \\
RRF $k$ & 60 & Cormack et al. (2009) \\
Seuil Jaccard & 0.85 & Manku et al. (2007) \\
$\varepsilon$ rerank & $10^{-9}$ & Stabilité numérique \\
Chunk size & 512 & BERT/MPNet \\
Chunk overlap & 64 & Lewis et al. (2020) \\
Top-K & 5 & Budget LLM \\
MIN\_PER\_DOC & 2 & Couverture \\
BM25 $k_1$ & 1.5 & Standard \\
BM25 $b$ & 0.75 & Standard \\
\bottomrule
\end{tabular}
\end{center}

\vspace{1cm}
\noindent Rapport rédigé par l'équipe PFE --- Mars 2026\\
Bettayeb Mohamed Aimen | Amina | Yanis | Loubna

\end{document}
